% ============================================================
%   字号、行距：不写单位，默认 pt
%   边距、间距、线宽：必须带单位
%   字重：\bfseries 为粗体，\mediumseries 为中等，\mdseries 为常规。
%   灰度：0 为黑色，1 为白色，数值越大颜色越浅。
% ============================================================

% ---------- 页面 ----------
\newcommand{\pageLeftMargin}{13mm}       % 左页边距
\newcommand{\pageRightMargin}{13mm}      % 右页边距
\newcommand{\pageTopMargin}{15mm}        % 上页边距
\newcommand{\pageBottomMargin}{12mm}     % 下页边距

% ---------- 文档大标题与副标题 ----------
\newcommand{\titleSize}{17}              % Itinerary / Resume / Booking 字号
\newcommand{\titleLeading}{19}           % 大标题行距
\newcommand{\titleWeight}{\bfseries}     % 大标题字重
\newcommand{\subtitleSize}{8}            % 签证副标题字号
\newcommand{\subtitleLeading}{9}         % 副标题行距
\newcommand{\subtitleWeight}{\mdseries}  % 副标题字重
\newcommand{\titleSubtitleGap}{5pt}      % 大标题与副标题之间的间距
\newcommand{\subtitleBodyGap}{10pt}      % 首个模块的上间距；保留 center 环境自带留白

% ---------- 模块标题 ----------
\newcommand{\sectionTitleSize}{11}       % PERSONAL INFORMATION 等模块标题字号
\newcommand{\sectionTitleLeading}{12}    % 模块标题行距
\newcommand{\sectionTitleWeight}{\bfseries}
\newcommand{\sectionTopGap}{20pt}        % 普通模块标题上方间距；首个模块用 subtitleBodyGap
\newcommand{\sectionRuleGap}{3pt}        % 模块标题文字与下方横线的间距
\newcommand{\sectionBottomGap}{8pt}      % 模块横线与其下方内容的间距
\newcommand{\sectionRuleWidth}{0.65pt}   % 模块标题下方横线的粗细

% ---------- 正文与字段 ----------
\newcommand{\bodySize}{10}               % 普通正文：简历个人信息、实习内容、其它信息
\newcommand{\bodyLeading}{13.7}          % 普通正文行距；不改变表格和小字字段的独立字号
\newcommand{\bodyWeight}{\mdseries}      % 普通正文字重
\newcommand{\fieldKeyWeight}{\mediumseries} % 个人信息、其它信息的键名：中等字重，介于常规和粗体之间
\newcommand{\fieldValueWeight}{\mdseries}% 字段值字重
\newcommand{\smallValueSize}{9}          % 行程摘要：字段值字号
\newcommand{\smallValueLeading}{11}      % 对应字段值行距
\newcommand{\summaryKeySize}{7.5}        % Trip Summary 左侧字段键字号
\newcommand{\summaryKeyWeight}{\bfseries}% Trip Summary 字段键字重，独立于个人信息
\newcommand{\summaryKeyValueGap}{10mm}   % Trip Summary 左侧标签与右侧内容的间距
\newcommand{\summaryRowGap}{4pt}         % Trip Summary 相邻两行的附加间距

% ---------- 个人信息：三份文档共用 ----------
% 所有字段键自动转为大写，不带冒号；每两个字段一行，宽字段独占一行。
\newcommand{\personalColumnGap}{8mm}     % 个人信息左右两组字段之间的间距
\newcommand{\personalLeftKeyWidth}{31mm}         % 左侧字段键列宽
\newcommand{\personalRightKeyWidth}{36mm}        % 右侧字段键列宽
\newcommand{\personalKeyValueGap}{2mm}           % 个人信息键与值之间的间距
\newcommand{\personalRowGap}{1.7mm}      % 个人信息每行之后的间距

% ---------- 表格：三份文档共用 ----------
\newcommand{\tableBodySize}{9.2}         % 表格正文字号，独立于普通正文
\newcommand{\tableBodyLeading}{11.4}     % 表格正文基础行距
\newcommand{\tableHeaderSize}{7.5}       % 表头字号
\newcommand{\tableHeaderLeading}{8.5}    % 表头行距
\newcommand{\tableHeaderWeight}{\bfseries}
\newcommand{\tableCellPadding}{4pt}      % 单元格左右各自的留白
\newcommand{\tableRowStretch}{1.4}       % 行高倍率；越大越疏朗
\newcommand{\tableRowRuleWidth}{0.35pt}  % 数据行之间的细横线；最后一行下方不画线
\newcommand{\tableHeaderRuleWidth}{0.8pt}% 表头下方横线的粗细

% ---------- Itinerary 专用尺寸 ----------
\newcommand{\itineraryDateWidth}{12mm}         % 行程表日期列；PLAN 自动占据剩余宽度
\newcommand{\itineraryLocationWidth}{25mm}     % 行程表地点列
\newcommand{\itineraryAccommodationWidth}{32mm}% 行程表住宿列
\newcommand{\itineraryTransportWidth}{21mm}    % 行程表交通列
\newcommand{\noteSize}{9}                      % Note 正文字号
\newcommand{\noteLeading}{11.5}                % Note 行距
\newcommand{\noteLeftMargin}{12pt}            % Note 项目符号列表缩进
\newcommand{\noteItemGap}{3pt}                 % Note 各条之间的附加间距

% ---------- Resume 专用尺寸 ----------
\newcommand{\resumeTableFirstWidth}{60mm}      % 教育/出国表格第一列宽度
\newcommand{\resumeTableSecondWidth}{76mm}     % 两张表格第二列宽度；日期列使用剩余宽度
\newcommand{\internshipCompanyWidth}{60mm}     % 实习公司名称列宽
\newcommand{\internshipColumnGap}{4mm}        % 公司、职位、日期之间的间距
\newcommand{\internshipCompanyWeight}{\bfseries}
\newcommand{\internshipRoleShape}{\itshape}   % 职位斜体；改为 \upshape 则使用直立体
\newcommand{\internshipDateSize}{9.5}         % 实习起止日期字号
\newcommand{\internshipDateLeading}{13}       % 实习起止日期行距
\newcommand{\internshipTopGap}{3mm}           % 每段实习上方的间距
\newcommand{\internshipOverviewGap}{0.8mm}    % 非空实习概述与标题的间距
\newcommand{\internshipListLeftMargin}{4.7mm} % 实习职责列表缩进
\newcommand{\internshipBulletGap}{2mm}        % 项目符号与文字的间距
\newcommand{\internshipListTopGap}{1.5mm}     % 职责列表上下留白
\newcommand{\internshipItemGap}{1.3mm}        % 三条实习职责之间的附加间距
\newcommand{\otherKeyWidth}{43mm}             % Other Information 字段键列宽
\newcommand{\otherKeyValueGap}{3mm}           % Other Information 键与值的间距
\newcommand{\otherRowGap}{1.5mm}              % Other Information 每行之后的间距

% ---------- Booking 专用尺寸 ----------
\newcommand{\bookingFlightWidth}{20mm}       % 航班号列；出发/到达平分剩余宽度
\newcommand{\bookingCityWidth}{22mm}         % 住宿城市列
\newcommand{\bookingDateWidth}{25mm}         % 第二列：入住和退房日期
\newcommand{\bookingStayWidth}{37mm}         % 住宿名称列
\newcommand{\bookingContactWidth}{29mm}      % 住宿联系电话列
\newcommand{\bookingPostalWidth}{18mm}       % 邮编列；地址使用剩余宽度

% ---------- 公共宏包、字体与颜色 ----------
% 使用 XeLaTeX；优先使用系统字体，缺失时自动使用 TeX 自带字体。
\usepackage[a4paper,left=\pageLeftMargin,right=\pageRightMargin,
  top=\pageTopMargin,bottom=\pageBottomMargin]{geometry}
\usepackage{lmodern} % 提供可缩放的后备字体，避免小字号替换警告
\usepackage{fontspec,xeCJK}
\usepackage{array,tabularx,longtable,ragged2e,etoolbox,needspace}
\usepackage[table]{xcolor}
\usepackage{booktabs,microtype,enumitem}
\IfFontExistsTF{Helvetica Neue}{%
  \setmainfont{Helvetica Neue}[UprightFont=*,BoldFont=* Bold,
    ItalicFont=* Italic,BoldItalicFont=* Bold Italic,
    FontFace={mb}{n}{Font=* Medium}]
}{%
  \setmainfont{texgyreheros-regular.otf}[BoldFont=texgyreheros-bold.otf,
    ItalicFont=texgyreheros-italic.otf,BoldItalicFont=texgyreheros-bolditalic.otf,
    FontFace={mb}{n}{Font=texgyreheros-regular.otf,FakeBold=0.5}]
}
% 使用字体自带的 Medium；后备字体缺少中等字重时仅轻微加粗。
\newcommand{\mediumseries}{\fontseries{mb}\selectfont}
\IfFontExistsTF{PingFang SC}{\setCJKmainfont{PingFang SC}}{\setCJKmainfont{FandolHei-Regular.otf}}
\definecolor{Text}{gray}{0.08}        % 正文与模块横线
\definecolor{Muted}{gray}{0.38}       % 副标题与小字字段键
\definecolor{Rule}{gray}{0.82}        % 表格行间细横线
\definecolor{HeaderRule}{gray}{0.47}  % 表头下方横线

% ---------- 公共排版实现：通常只需修改上方参数 ----------
\color{Text}
\setlength{\parindent}{0pt}
\setlength{\parskip}{0pt}
\setlength{\tabcolsep}{0pt}
\setlength{\arrayrulewidth}{\tableRowRuleWidth}
\pagestyle{empty}
\newcommand{\bodyfont}{\fontsize{\bodySize}{\bodyLeading}\selectfont\bodyWeight}
\newcommand{\smallvaluefont}{\fontsize{\smallValueSize}{\smallValueLeading}\selectfont\fieldValueWeight}
\newcommand{\summarykeyfont}{\fontsize{\summaryKeySize}{\smallValueLeading}\selectfont\summaryKeyWeight\color{Muted}}
\newcommand{\visaCategory}[1]{\def\visaCategoryValue{#1}}

% 三份文档共用标题；签证类型来自各自 info.tex 中的 \visaCategory{...}。
\newcommand{\documenttitle}[1]{%
  \begingroup\bodyfont\RaggedRight
  \begin{center}
    {\fontsize{\titleSize}{\titleLeading}\selectfont\titleWeight #1\par}%
    \vspace{\titleSubtitleGap}%
    {\fontsize{\subtitleSize}{\subtitleLeading}\selectfont\subtitleWeight\color{Muted}\visaCategoryValue\ Visa Application\par}%
  \end{center}
  \endgroup}

% 可选参数覆盖当前模块上间距；首个模块使用 \subtitleBodyGap。
\newcommand{\sectionhead}[2][\sectionTopGap]{%
  \par\addvspace{#1}\Needspace{5\baselineskip}%
  {\fontsize{\sectionTitleSize}{\sectionTitleLeading}\selectfont\sectionTitleWeight\MakeUppercase{#2}\par}%
  \nobreak\vspace{\sectionRuleGap}{\color{Text}\hrule height \sectionRuleWidth}%
  \nobreak\vspace{\sectionBottomGap}\nointerlineskip}

% 个人信息的数据收集和布局由三份文档共用。
\newcommand{\personalrows}{}
\newcommand{\personalrow}[2]{\appto\personalrows{\renderpersonalfield{#1}{#2}}}
\newcommand{\personalwide}[2]{\appto\personalrows{\renderpersonalwide{#1}{#2}}}
\newsavebox{\personalleftbox}
\newif\ifpersonalpending
\newlength{\personalleftlabelwidth}
\newlength{\personalrightlabelwidth}
\setlength{\personalleftlabelwidth}{\personalLeftKeyWidth}
\setlength{\personalrightlabelwidth}{\personalRightKeyWidth}
\newcommand{\personalkeyvalue}[3]{%
  \begin{minipage}[t]{#1}%
    \RaggedRight\strut{\fieldKeyWeight\MakeUppercase{#2}}\strut
  \end{minipage}\hspace{\personalKeyValueGap}%
  \begin{minipage}[t]{\dimexpr\linewidth-#1-\personalKeyValueGap\relax}%
    \RaggedRight\fieldValueWeight\strut #3\strut
  \end{minipage}}
\newcommand{\personalfield}[3]{%
  \begin{minipage}[t]{\dimexpr(\linewidth-\personalColumnGap)/2\relax}%
    \personalkeyvalue{#3}{#1}{#2}%
  \end{minipage}}
\newcommand{\renderpersonalfield}[2]{%
  \ifpersonalpending
    \noindent\usebox{\personalleftbox}\hfill\personalfield{#1}{#2}{\personalrightlabelwidth}\par\vspace{\personalRowGap}%
    \personalpendingfalse
  \else
    \sbox{\personalleftbox}{\personalfield{#1}{#2}{\personalleftlabelwidth}}%
    \personalpendingtrue
  \fi}
\newcommand{\flushpersonalrow}{%
  \ifpersonalpending
    \noindent\usebox{\personalleftbox}\par\vspace{\personalRowGap}%
    \personalpendingfalse
  \fi}
\newcommand{\renderpersonalwide}[2]{%
  \flushpersonalrow
  \noindent\begin{minipage}[t]{\linewidth}%
    \personalkeyvalue{\personalleftlabelwidth}{#1}{#2}%
  \end{minipage}\par\vspace{\personalRowGap}}
\newcommand{\personalinfo}{%
  \begingroup\bodyfont\RaggedRight
  \personalpendingfalse
  \personalrows\flushpersonalrow
  \endgroup}

% 表格均左对齐；L 为固定宽度列，Y 自动占据剩余宽度。
\newcolumntype{L}[1]{>{\RaggedRight\arraybackslash}m{#1}}
\newcolumntype{Y}{>{\RaggedRight\arraybackslash}X}
\newcommand{\tablestyle}{%
  \fontsize{\tableBodySize}{\tableBodyLeading}\selectfont\bodyWeight
  \setlength{\tabcolsep}{\tableCellPadding}%
  \setlength{\arrayrulewidth}{\tableRowRuleWidth}%
  \renewcommand{\arraystretch}{\tableRowStretch}%
  \arrayrulecolor{Rule}}
\newcommand{\tableheadrule}{\arrayrulecolor{HeaderRule}\specialrule{\tableHeaderRuleWidth}{0pt}{0pt}\arrayrulecolor{Rule}}
\newcommand{\headcell}[1]{{\fontsize{\tableHeaderSize}{\tableHeaderLeading}\selectfont\tableHeaderWeight\color{Text}\MakeUppercase{#1}}}
